% -------------
\section{LLM Prompt for Domain Validation}
\label{sec:appendix_domain_validation}

To identify core domains, we consulted a panel of four frontier models: \texttt{GPT-4o},
\texttt{Claude 3.5 Sonnet}, \texttt{Gemini 2.5 Pro}, and \texttt{DeepSeek V3}.
Each was prompted individually to identify cornerstone engineering domains.
We finalized the 15 core domains via majority voting.
The domain identification prompt is provided below.

\begin{tcolorbox}[
    colback=gray!5,
    colframe=gray!40,
    boxrule=0.4pt,
    arc=2pt,
    left=6pt,
    right=6pt,
    top=6pt,
    bottom=6pt,
    breakable
]
\ttfamily\small
You are an expert academic and senior curriculum designer at a top-tier engineering university (like MIT or Caltech). You have decades of experience structuring undergraduate programs for ABET accreditation. \\

Your task is to identify and rank the most fundamental, cornerstone
domains for the five primary branches of engineering. The goal is to identify the non-overlapping, core subject areas that are prerequisites for almost all other advanced topics. \\

For example, for Chemical Engineering, ``Reaction Kinetics" would be a fundamental domain, while a niche topic like ``Polymer Rheology" would be a sub-specialization, not a cornerstone. \\

Please provide a ranked list of the top 5--7 most fundamental domains for each of the following five branches: \\
1. Chemical Engineering \\
2. Electrical Engineering \\
3. Mechanical Engineering \\
4. Civil Engineering \\
5. Industrial Engineering \\

Return your answer \textbf{only} in a strict JSON format as given below.

\vspace{8pt}
\noindent\bgroup\ttfamily
\texttt{\char`\{}

\vspace{8pt}
\hspace*{1em}``chemical\_engineering": [\\
\hspace*{2em}``Rank 1 Domain",\\
\hspace*{2em}``Rank 2 Domain",\\
\hspace*{2em}``Rank 3 Domain",\\
\hspace*{2em}``Rank 4 Domain",\\
\hspace*{2em}``Rank 5 Domain"\\
\hspace*{1em}],\\

\vspace{2pt}
\hspace*{1em}``electrical\_engineering": [\\
\hspace*{2em}``Rank 1 Domain",\\
\hspace*{2em}``Rank 2 Domain",\\
\hspace*{2em}``Rank 3 Domain",\\
\hspace*{2em}``Rank 4 Domain",\\
\hspace*{2em}``Rank 5 Domain"\\
\hspace*{1em}],\\

\vspace{2pt}
\hspace*{1em}``mechanical\_engineering": [\\
\hspace*{2em}``Rank 1 Domain",\\
\hspace*{2em}``Rank 2 Domain",\\
\hspace*{2em}``Rank 3 Domain",\\
\hspace*{2em}``Rank 4 Domain",\\
\hspace*{2em}``Rank 5 Domain"\\
\hspace*{1em}],\\

\vspace{2pt}
\hspace*{1em}``civil\_engineering": [\\
\hspace*{2em}``Rank 1 Domain",\\
\hspace*{2em}``Rank 2 Domain",\\
\hspace*{2em}``Rank 3 Domain",\\
\hspace*{2em}``Rank 4 Domain",\\
\hspace*{2em}``Rank 5 Domain"\\
\hspace*{1em}],\\

\vspace{2pt}
\hspace*{1em}``industrial\_engineering": [\\
\hspace*{2em}``Rank 1 Domain",\\
\hspace*{2em}``Rank 2 Domain",\\
\hspace*{2em}``Rank 3 Domain",\\
\hspace*{2em}``Rank 4 Domain",\\
\hspace*{2em}``Rank 5 Domain"\\
\hspace*{1em}]\\

\vspace{2pt}
\texttt{\char`\}}
\egroup
\vspace{8pt}

Do not include any preamble, conversational text, explanations, or markdown formatting around the JSON block.
\end{tcolorbox}

% -------------
\section{LLM Prompt for Area Validation}
\label{sec:appendix_area_validation}

Following the domain validation procedure, we utilized the same model panel and aggregated the
pedagogical areas via majority voting.
The prompt used for area identification is provided below.

\begin{tcolorbox}[
    colback=gray!5,
    colframe=gray!40,
    boxrule=0.4pt,
    arc=2pt,
    left=6pt,
    right=6pt,
    top=6pt,
    bottom=6pt,
    breakable
]
\ttfamily\small
You are an expert academic and senior curriculum designer at a top-tier engineering university (like MIT or Caltech). You have decades of experience structuring undergraduate programs for ABET accreditation. \\

Your task is to identify and list the most fundamental, non-overlapping ``areas" or
``sub-topics" within the provided engineering domain. \\

A ``fundamental area" is a major pedagogical unit or chapter within a core course on this domain. \textbf{The areas you list should be distinct pillars of the subject.} For example, if the domain is
``Thermodynamics," a fundamental area would be ``The First Law" or
``Properties of Pure Fluids." A concept like ``Enthalpy" would be
too specific. Similarly, within a ``Circuit Analysis" domain, ``Kirchhoff's Voltage Law" and ``Kirchhoff's Current Law" are foundational concepts often taught together and would be too granular to list separately. \\

Please provide a list of the top 3-5 most fundamental areas for
the following engineering domain:
\texttt{[DOMAIN\_NAME\_HERE]} \\

Return your answer \textbf{only} in a strict JSON format as given below.

\vspace{8pt}
\noindent\bgroup\ttfamily
\texttt{\char`\{}

\vspace{4pt}
\hspace*{1em}``domain": ``[DOMAIN\_NAME\_HERE]",\\
\hspace*{1em}``areas": [\\
\hspace*{2em}``Area 1",\\
\hspace*{2em}``Area 2",\\
\hspace*{2em}``Area 3",\\
\hspace*{2em}``Area 4"\\
\hspace*{1em}]\\
\texttt{\char`\}}
\egroup
\vspace{8pt}

Do not include any preamble, conversational text, explanations, or markdown formatting around the JSON block.
\end{tcolorbox}

% -------------
\section{Foundational Textbooks and Authoritative Data Sources}
\label{sec:appendix_books}
\label{sec:appendix_data_sources}

The core engineering principles, problem typologies, and realistic constraints for the templates
developed by us were sourced from the authoritative textbooks listed
in~\autoref{tab:foundational_textbooks}.
These texts, which are standard in top-tier engineering curricula, served as the primary knowledge
base to ensure \ourdataset is pedagogically sound and grounded in the core engineering curriculum.
They were also used to validate our selection of fundamental domains and areas.

\begin{table}[t]
\small
\centering
\renewcommand{\arraystretch}{1.1}
\setlength{\tabcolsep}{6pt}
\begin{tabular}{p{0.35\linewidth} p{0.55\linewidth}}
\toprule
\rowcolor{gray!10}\textbf{Domain} & \textbf{Textbook} \\
\midrule
Reaction Kinetics & Elements of Chemical Reaction Engineering by H. Scott Fogler \\
Thermodynamics & Introduction to Chemical Engineering Thermodynamics by J. M. Smith, H. C. Van Ness,
and M. M. Abbott \\
Transport Phenomena & Transport Phenomena by R. Byron Bird, Warren E. Stewart, and Edwin N.
Lightfoot \\
\addlinespace[2pt]
Electromagnetics and Waves & Fundamentals of Applied Electromagnetics by F. T. Ulaby and Umberto
Ravaioli \\
Digital Communications & Digital Communications by John G. Proakis and Masoud Salehi \\
Signals and Systems & Discrete-Time Signal Processing by Alan V. Oppenheim and Ronald W. Schafer \\
\addlinespace[2pt]
Mechanics of Materials & Mechanics of Materials by Ferdinand P. Beer, E. Russell Johnston, Jr., and
John T. DeWolf \\
Fluid Mechanics & Fundamentals of Fluid Mechanics by Bruce R. Munson, Donald F. Young, and Theodore
H. Okiishi \\
Vibrations and Acoustics & Mechanical Vibrations by Singiresu S. Rao \\
\addlinespace[2pt]
Structural Analysis & Structural Analysis by R. C. Hibbeler \\
Geotechnical Engineering & Principles of Geotechnical Engineering by Braja M. Das and Khaled
Sobhan \\
Water Resources & Open Channel Hydraulics by Terry W. Sturm \newline
Applied Hydrology by Ven Te Chow, David R. Maidment, and Larry W. Mays \\
\addlinespace[2pt]
Production and Inventory & Production and Operations Analysis by Steven Nahmias and Tava Lennon
Olsen \\
Quality and Reliability Control & Introduction to Statistical Quality Control by Douglas C.
Montgomery \\
Stochastic Operations & Introduction to Probability Models by Sheldon M. Ross \newline
Introduction to Operations Research by Frederick S. Hillier and Gerald J. Lieberman \\
\bottomrule
\end{tabular}
\caption{Foundational textbooks used for selecting domain-specific problems.}
\label{tab:foundational_textbooks}
\end{table}

To ensure the physical and engineering realism of \ourdataset, we built our domain-aware
parameterization on authoritative data.
We extracted extensive lists of physical constants and material properties from the standard
handbooks and data sources listed in~\autoref{tab:authoritative_sources}, and each constants table
in the code records the provenance of its values.

\begin{table}[t]
\small
\centering
\renewcommand{\arraystretch}{1.1}
\setlength{\tabcolsep}{6pt}
\begin{tabular}{p{0.35\linewidth} p{0.50\linewidth}}
\toprule
\rowcolor{gray!10}\textbf{Engineering Branch} & \textbf{Authoritative Data Sources} \\
\midrule
Chemical Engineering &
Perry's Chemical Engineers' Handbook \newline
NIST Chemistry WebBook \newline
CRC Handbook of Chemistry and Physics \newline
NIST-JANAF Thermochemical Tables \newline
NASA Technical Report R-132 \\
\addlinespace[2pt]
Electrical Engineering &
IEEE 100: The Authoritative Dictionary of IEEE Standards Terms \newline
Standard Handbook for Electrical Engineers \newline
ART-DEIT Database \newline
CODATA 2022 Recommended Values \newline
refractiveindex.info Database \\
\addlinespace[2pt]
Mechanical Engineering &
ASM Handbook, Volume 1 \newline
Marks' Standard Handbook for Mechanical Engineers \newline
Shigley's Mechanical Engineering Design \newline
MIL-HDBK-5J \newline
USDA Wood Handbook \newline
NIST Chemistry WebBook \newline
PubChem \newline
NIST SP 811 \\
\addlinespace[2pt]
Civil Engineering &
NAVFAC DM-7.01 and DM-7.02 \newline
FHWA HDS-4 and HEC-22 \newline
NRCS TR-55 \newline
USGS Water-Supply Paper 2339 \newline
AISC Shapes Database v16.0 \newline
NIST Chemistry WebBook \newline
CODATA 2022 Recommended Values \newline
NIST SP 811 \\
\addlinespace[2pt]
Industrial Engineering &
MIL-STD-105E \newline
NIST/SEMATECH e-Handbook of Statistical Methods \\
\bottomrule
\end{tabular}
\caption{Primary data sources referenced for engineering parameterization across branches.}
\label{tab:authoritative_sources}
\end{table}

% -------------
\section{LLM Prompt for Pedagogical Scoring}
\label{sec:appendix_significance_scoring}

The following prompt was used to query LLMs (described in~\autoref{sec:appendix_domain_validation})
to evaluate the pedagogical significance of each sub-topic within its parent domain.

\begin{tcolorbox}[
    colback=gray!5,
    colframe=gray!30,
    boxrule=0.4pt,
    arc=2pt,
    left=6pt,
    right=6pt,
    top=6pt,
    bottom=6pt,
    breakable
]
\ttfamily\small
You are an expert academic and senior curriculum designer at a top-tier engineering university with decades of experience structuring undergraduate programs for ABET accreditation. \\

Your task is to evaluate the pedagogical significance of a specific engineering ``area" within a broader ``domain" on a scale of 1 to 5, based on the rubric below. \\

\textbf{Rubric:}
\par\medskip\noindent
\textbf{Score 5 (Cornerstone):} A
foundational, prerequisite topic for nearly all other concepts in the domain.
\par\medskip\noindent
\textbf{Score 3 (Core Concept):} A standard, important topic that builds upon cornerstone principles but is not necessarily a universal prerequisite.
\par\medskip\noindent
\textbf{Score 1 (Specialized Application):} A focused application or an
integrative topic taught after cornerstones and core concepts are mastered.
\par\medskip

Please evaluate the following area:
\par\medskip\noindent
\textbf{Domain:} [DOMAIN\_NAME\_HERE]
\par\noindent
\textbf{Area:} [AREA\_NAME\_HERE]
\par\bigskip

\noindent\bgroup\ttfamily
\texttt{\char`\{}

\par\medskip
\hspace*{1em}``area": ``[AREA\_NAME\_HERE]",\\
\hspace*{1em}``pedagogical\_significance\_score": <integer, a score from 1 to 5>,\\
\hspace*{1em}``justification": <string, a one-sentence explanation for the score>.\\
\vspace{2pt}
\texttt{\char`\}}
\egroup
\par\bigskip

Do not include any preamble, conversational text, explanations, or markdown formatting around the JSON block.
\end{tcolorbox}
